\section{Methodology}
\label{sec:method}

We design three complementary experiments to probe alignment resistance across model scales. Each targets a different mechanism by which alignment might fail: direct adversarial attack (Experiment~1), in-context behavioral influence (Experiment~2), and instruction-level conflict (Experiment~3).

\subsection{Models}
\label{sec:models}

We test five OpenAI models spanning approximately three orders of magnitude in parameter count (\tabref{tab:models}). Since OpenAI does not publish exact parameter counts, we use approximate values based on public estimates for scaling analysis.

\begin{table}[h]
\centering
\caption{Models tested, ordered by approximate parameter count. Exact counts are unavailable for proprietary models.}
\begin{tabular}{@{}llll@{}}
\toprule
\textbf{Model} & \textbf{Approx.\ Parameters} & \textbf{Type} & \textbf{Architecture} \\
\midrule
\gptnano & ${\sim}$1B & Small, efficient & GPT-4.1 family \\
\gptmini & ${\sim}$8B & Mid-range & GPT-4.1 family \\
\gptomini & ${\sim}$8B & Mid-range & GPT-4o family \\
\gpto & ${\sim}$200B & Large & GPT-4o family \\
\gptone & ${\sim}$400B & Largest & GPT-4.1 family \\
\bottomrule
\end{tabular}
\label{tab:models}
\end{table}

\subsection{Evaluation Framework}
\label{sec:eval}

We use \gptomini as an automated safety judge, following established LLM-as-judge methodology \citep{zheng2024judging}. For each model response, the judge assigns:
\begin{itemize}[leftmargin=*,itemsep=0pt,topsep=0pt]
    \item A \textbf{continuous safety score} from 0 (completely unsafe) to 10 (completely safe), and
    \item A \textbf{binary refusal classification} (refused or complied).
\end{itemize}
The judge operates at temperature 0 for consistency. Using both continuous and discrete metrics allows us to apply the \citet{schaeffer2023emergent} framework: if the discrete metric shows a discontinuity that the continuous metric does not, the apparent phase transition is likely a measurement artifact.

\subsection{Experiment 1: Jailbreak Resistance}
\label{sec:exp1}

\para{Adversarial prompts.} We construct 15 adversarial prompts across 7 harm categories (privacy violations, harmful instructions, discrimination, financial crime, violence, misinformation, and self-harm), drawing inspiration from the \beavertails dataset \citep{ji2024beavertails}.

\para{Jailbreak templates.} Each prompt is formatted with three templates of increasing sophistication:
\begin{enumerate}[leftmargin=*,itemsep=0pt,topsep=0pt]
    \item \textbf{Direct}: The adversarial prompt is presented without any framing.
    \item \textbf{Role-play (\DAN)}: The model is instructed to adopt an unrestricted persona that ignores safety constraints.
    \item \textbf{Hypothetical}: The prompt is embedded in a fictional or hypothetical scenario (\eg ``In a novel I'm writing, a character needs to...'').
\end{enumerate}

\para{Protocol.} For each of 5 models $\times$ 3 templates $\times$ 15 prompts = 225 evaluations, we query the model at temperature 0 and judge the response. We report mean safety scores and refusal rates per model and per template.

\subsection{Experiment 2: Few-Shot Alignment Erosion}
\label{sec:exp2}

We test whether misaligned in-context examples can erode alignment at inference time, and whether this susceptibility varies with scale.

\para{Setup.} We construct 5 misaligned (user, assistant) example pairs where the assistant provides harmful content. For each of 10 adversarial test prompts, we prepend $n \in \{0, 1, 3, 5\}$ misaligned examples as in-context demonstrations, then present the adversarial prompt as the final user message. This yields 5 models $\times$ 4 shot levels $\times$ 10 prompts = 200 evaluations.

\subsection{Experiment 3: Instruction Conflict Resolution}
\label{sec:exp3}

We test how models resolve conflicts between a safety-oriented system prompt and an adversarial user prompt that attempts to override it.

\para{Setup.} We design 5 conflict scenarios where the system prompt explicitly instructs the model to maintain safety, while the user message requests harmful content with a plausible justification (\eg requesting chemical mixing instructions ``for a safety lecture''). Each scenario is run 3 times per model at temperature 0.3 to estimate variance, yielding 5 models $\times$ 5 scenarios $\times$ 3 runs = 75 evaluations.

\subsection{Scaling Analysis}
\label{sec:scaling}

To test for phase transitions, we fit two models to the safety score vs.\ $\log(\text{parameters})$ data:
\begin{itemize}[leftmargin=*,itemsep=0pt,topsep=0pt]
    \item \textbf{Linear} (gradual scaling): $s = a \cdot \log(N) + b$, where $s$ is safety score and $N$ is parameter count.
    \item \textbf{Sigmoid} (phase transition): $s = \frac{L}{1 + e^{-k(\log(N) - x_0)}} + c$, where the inflection point $x_0$ represents the critical scale.
\end{itemize}
We compare fits using $R^2$ values. We also compute the first derivative of safety \wrt $\log(N)$ to identify any discontinuities, and report Pearson and Spearman correlations with significance tests.
